\documentclass[11pt,a4paper]{article}

\usepackage[utf8]{inputenc}
\usepackage[T1]{fontenc}
\usepackage[brazil]{babel}
\usepackage{lmodern}
\usepackage[margin=2.3cm]{geometry}
\usepackage{booktabs}
\usepackage{longtable}
\usepackage{array}
\usepackage{xcolor}
\usepackage{enumitem}
\usepackage{listings}
\usepackage{microtype}
\usepackage{tikz}
\usetikzlibrary{shapes.geometric, arrows.meta, positioning, fit, backgrounds, calc}
\usepackage[hidelinks]{hyperref}
\usepackage{caption}

\definecolor{verde}{RGB}{22,128,60}
\definecolor{ambar}{RGB}{176,116,0}
\definecolor{vermelho}{RGB}{176,26,26}
\definecolor{azul}{RGB}{37,70,160}
\definecolor{cinza}{RGB}{95,95,95}
\definecolor{fundo}{RGB}{247,247,247}
\definecolor{fundoazul}{RGB}{234,240,252}
\definecolor{fundoverde}{RGB}{233,245,237}
\definecolor{fundoambar}{RGB}{253,245,230}

\newcommand{\Verde}[1]{\textcolor{verde}{\textbf{#1}}}
\newcommand{\Ambar}[1]{\textcolor{ambar}{\textbf{#1}}}
\newcommand{\Vermelho}[1]{\textcolor{vermelho}{\textbf{#1}}}
\newcommand{\cod}[1]{\texttt{\small #1}}

\lstset{
  basicstyle=\ttfamily\footnotesize, backgroundcolor=\color{fundo}, frame=single,
  rulecolor=\color{cinza!40}, breaklines=true, columns=fullflexible,
  showstringspaces=false, xleftmargin=4pt, framexleftmargin=4pt,
  literate={á}{{\'a}}1 {é}{{\'e}}1 {í}{{\'i}}1 {ó}{{\'o}}1 {ú}{{\'u}}1
           {ã}{{\~a}}1 {õ}{{\~o}}1 {â}{{\^a}}1 {ê}{{\^e}}1 {ô}{{\^o}}1
           {ç}{{\c c}}1 {à}{{\`a}}1 {É}{{\'E}}1 {Á}{{\'A}}1 {Ç}{{\c C}}1
           {º}{{\textordmasculine}}1 {ª}{{\textordfeminine}}1 {–}{{--}}1 {—}{{---}}1
}

\tikzset{
  etapa/.style={rectangle, rounded corners=3pt, draw=azul!60, fill=fundoazul, text width=2.6cm,
                align=center, font=\scriptsize, minimum height=0.85cm, inner sep=3pt},
  portao/.style={rectangle, draw=cinza!70, fill=fundo, text width=2.6cm, align=center,
                 font=\scriptsize\itshape, minimum height=0.75cm, inner sep=3pt},
  loopbox/.style={rectangle, rounded corners=5pt, draw=ambar!80, dashed, thick, inner sep=7pt},
  seta/.style={-{Stealth[length=2.2mm]}, draw=cinza!85},
  setavolta/.style={-{Stealth[length=2.2mm]}, draw=ambar, dashed, thick},
}

\title{\textbf{Gerador de Missões Unooria}\\[0.35em]
\large Guia de funcionamento: o que a esteira faz, por que faz assim,\\ e como escrever uma entrada que produz uma boa missão}
\author{Documentação técnica --- repositório \texttt{gerador-missoes}}
\date{19 de setembro de 2026}

\begin{document}
\maketitle

\begin{abstract}
\noindent
O gerador transforma uma ideia de missão em um \textbf{PRD pronto para codificação} e um
\textbf{one-pager que o time pedagógico aprova}, passando por 50 agentes organizados em três
etapas e seis ciclos de correção automática. Este guia é para quem vai \emph{usar} a
ferramenta: explica o que colocar na entrada, o que cada etapa faz, como ler o diagnóstico
quando algo trava, e por que o sistema foi desenhado desta forma e não como um único
prompt grande. A Seção~\ref{sec:entrada} --- com exemplos de briefing bom e ruim --- é a que
mais muda o resultado na prática.
\end{abstract}

\tableofcontents
\newpage

\section{O que a esteira faz}

Você escreve um briefing de meia página sobre a missão que quer. A esteira devolve:

\begin{table}[h]
\centering
\begin{tabular}{p{4.6cm}p{10.4cm}}
\toprule
\textbf{Artefato} & \textbf{Para quem é} \\
\midrule
\cod{one-pager-pedagogico.md} (+ \cod{.docx}) & Time pedagógico, coordenação e a professora consultora. Linguagem de sala de aula, zero jargão. É o documento de aprovação. \\
\cod{PRD.md} & O agente (ou a pessoa) que vai codificar. 25 seções, com contratos de dados, eventos e arquitetura. \\
\cod{diagnostico-execucao.md} & Você. Diz o que passou, o que passou com ressalva, o que travou e de quem é cada conserto. \\
\cod{design/09-diagramas.html} & Quem quiser \emph{ver} a missão: jornada e fluxo em diagramas. \\
\cod{design/08-prototipo.html} & Revisão humana: as telas navegáveis, abre com duplo clique. \\
\cod{aprendizado.md} & O processo: o que esta execução ensinou e o que propor de mudança nas rubricas. \\
\bottomrule
\end{tabular}
\end{table}

\noindent
Mais os documentos intermediários: spec da missão, jornada, fluxo, copy, especificação de
telas, camada de acessibilidade, handoff de dev e os relatórios de cada portão.

\subsection{A regra que muda tudo: a esteira não desiste}

Um portão reprovado \textbf{não} interrompe a execução. Todos os documentos são produzidos
até o fim; o que deriva de uma etapa reprovada sai marcado \cod{[PROVISÓRIO]}, e o
diagnóstico final diz o que é firme e o que é rascunho.

Isso é deliberado. Na versão anterior, a primeira reprovação parava tudo: não saía PRD, não
saía one-pager, não saía relatório --- e quem precisava entender o erro ficava sem material
para analisar. Um sistema que falha sem deixar evidência transfere para a pessoa o trabalho
de reconstruir o que aconteceu.

\section{Como escrever a entrada}
\label{sec:entrada}

Esta é a seção com maior efeito sobre o resultado. A esteira é boa em \emph{estruturar},
\emph{auditar} e \emph{formalizar}. Ela é ruim em adivinhar o que você não disse --- e, por
desenho, prefere marcar uma lacuna a inventar conteúdo.

\subsection{O que o primeiro agente precisa}

O agente 01 (Núcleo Narrativo) produz quatro coisas a partir do seu briefing: a
\textbf{descoberta central}, a \textbf{narrativa}, o \textbf{mecanismo das seis etapas} e a
\textbf{identificação} (ano, bimestre, ambiente). Ele precisa de matéria-prima para as
quatro.

\begin{center}
\begin{tikzpicture}[node distance=3.05cm]
\node[etapa] (a) {\textbf{Desafiar}\\ a tarefa concreta};
\node[etapa, right=of a] (b) {\textbf{Navegar}\\ explorar sem saber};
\node[etapa, right=of b] (c) {\textbf{Decidir}\\ a escolha que importa};
\node[etapa, below=0.9cm of c] (d) {\textbf{Resultado}\\ o ``ahá''};
\node[etapa, left=of d] (e) {\textbf{Explicar}\\ por que aconteceu};
\node[etapa, left=of e] (f) {\textbf{Refletir}\\ aplicar em cenário novo};
\draw[seta] (a) -- (b); \draw[seta] (b) -- (c); \draw[seta] (c) -- (d);
\draw[seta] (d) -- (e); \draw[seta] (e) -- (f);
\node[font=\scriptsize\itshape, text width=6cm, align=center, below=0.35cm of e, text=cinza]
  {O estudante que vive o Navegar \textbf{não pode saber} o que o Resultado vai revelar. Esse é o teste mais duro da esteira.};
\end{tikzpicture}
\captionof{figure}{O framework de seis etapas que toda missão segue}
\end{center}

\subsection{Um briefing que funciona}

\begin{lstlisting}
Missao para o 5o ano do Fundamental, 3o bimestre, no ambiente do YouTube (unitube).

Tema do bimestre (material Entre Redes anexado): como a recomendacao de
conteudo funciona e por que duas pessoas veem coisas diferentes.

O que quero que a crianca descubra: que o que ela ve no feed nao e "o que
existe", e sim o que o sistema aprendeu que ela costuma assistir.

Ideia de mecanismo: ela monta a playlist de um personagem a partir de um banco
de videos. Depois ve o feed que aquele personagem passa a receber. Dois colegas
com playlists diferentes recebem feeds diferentes, e isso fica visivel na tela.

Restricoes: 50 minutos, aula unica, dois estudantes por tablet.
A professora precisa conseguir pausar a turma antes da revelacao.

Em aberto: nao sei se o banco deve ter 8 ou 12 videos.
\end{lstlisting}

\noindent
Por que funciona:

\begin{enumerate}[leftmargin=*,itemsep=2pt]
  \item \textbf{Ano e bimestre explícitos.} O identificador da missão codifica os dois
  (Seção~\ref{sec:id}); sem eles o primeiro portão determinístico já reprova.
  \item \textbf{A descoberta está escrita como descoberta}, não como objetivo de aula. Compare:
  ``a criança vai aprender sobre algoritmos de recomendação'' (objetivo) contra ``o que ela vê
  não é o que existe'' (descoberta).
  \item \textbf{Há um mecanismo, não só uma intenção.} ``Monta a playlist, vê o feed que
  resulta, compara com o colega'' é uma cadeia de ações. ``Trabalhar pensamento crítico sobre
  algoritmos'' não é.
  \item \textbf{As restrições de sala aparecem.} Tempo, dispositivo compartilhado e controle
  da professora mudam o desenho, e a esteira audita os três.
  \item \textbf{A dúvida está declarada.} Dizer ``não sei se 8 ou 12'' faz a esteira decidir e
  registrar a decisão. Omitir faria ela decidir sem você saber.
\end{enumerate}

\subsection{Briefings que produzem missão ruim}

\begin{center}
\begin{longtable}{p{7.1cm}p{7.9cm}}
\toprule
\textbf{Entrada problemática} & \textbf{O que acontece} \\
\midrule
\endfirsthead
\multicolumn{2}{l}{\small\itshape continuação} \\
\toprule
\textbf{Entrada problemática} & \textbf{O que acontece} \\
\midrule
\endhead
``Missão sobre fake news para o Fundamental.'' &
Sem ano, sem bimestre, sem ambiente, sem mecanismo. A esteira inventa os quatro, e o que ela
inventa é o mais genérico possível. \Vermelho{Pior caso de uso da ferramenta.} \\
\addlinespace
``O estudante navega por um feed com posts falsos e descobre que foram manipulados.'' &
O briefing \textbf{entrega a virada}. A esteira tem um portão dedicado a isso (agente 02) e vai
reprovar, mas o estrago já aconteceu: a ideia nasceu sem segredo, e refazer o núcleo custa rodadas. \\
\addlinespace
``Uma missão que desenvolva pensamento crítico e cidadania digital.'' &
Intenção sem mecanismo. O gate estrutural tem um critério só para isso: se você remover a
frase, a missão continua igual, então a frase não descreve nada construível. \\
\addlinespace
``O estudante avalia 15 perfis e classifica cada um como confiável ou não.'' &
Passa no teste de decisão simples e reprova no volume: 15 decisões do mesmo tipo excede o teto
da faixa. Vira \Ambar{ressalva} ou reprovação no agente 09. \\
\addlinespace
``Fazer no Roblox, que as crianças amam.'' &
Ambiente fora da lista fechada (Instagram, YouTube, ChatGPT, TikTok, WhatsApp). Reprovação
imediata --- não existe seção de design system para construir. \\
\addlinespace
``A criança conversa com a IA e reflete sobre o que aprendeu.'' &
Reflexão genérica. O critério 1.6 pede uma \emph{provocação concreta}: uma situação nova em que
ela aplica o que viveu. O teste é direto: alguém que não fez a missão consegue responder? \\
\addlinespace
Briefing de oito páginas com todas as telas já desenhadas. &
A esteira vira transcritora. Os portões detectam ``eco'' --- quando um agente repete a entrada
em vez de construir --- mas você terá pago 50 chamadas para formatar o que já tinha. \\
\bottomrule
\end{longtable}
\end{center}

\subsection{Anexos valem mais que adjetivos}

Anexe o material do bimestre (o \emph{Entre Redes} em \cod{.docx}). Ele é a fonte primária de
código curricular e de contexto: a esteira lê os códigos da BNCC dali e os confere contra o
anexo oficial do MEC~\cite{cne2022}. Dois parágrafos de material real valem mais que uma
página de descrição.

\section{Por que 50 agentes, e não um prompt grande}

A pergunta aparece sempre, e a resposta não é ``porque especialização é boa''.

\subsection{Densidade de instrução degrada o cumprimento}

A etapa de Produto sozinha carrega cerca de 74 mil \emph{tokens} de instrução e 51 critérios
tabelados. No \emph{benchmark} IFScale, que mede exatamente isso, os melhores modelos de
fronteira chegam a 68\% de acurácia quando precisam seguir 500 instruções ao mesmo
tempo~\cite{ifscale2025}, com viés sistemático em favor das instruções apresentadas primeiro.
Liu et al.~\cite{liu2024} documentam o efeito irmão em contexto longo: o que está no meio é
usado pior que o que está nas pontas.

Há evidência interna, e ela é mais convincente que a externa. O prompt do agente 09 abre
justificando a própria existência: \textit{``esta avaliação existe porque a rubrica original
errava sistematicamente aqui: os critérios de faixa etária competiam por atenção com $\sim$30
critérios de outra natureza''}. A versão agregada falhava de modo reprodutível; a decomposta
corrigiu.

\subsection{O critério de decomposição}

Não se divide por etapa cronológica. Divide-se por \textbf{decisão encapsulada}, no sentido
clássico de Parnas~\cite{parnas1972}: cada agente responde uma pergunta que pode ser
respondida --- e errada --- independentemente das outras, e emite um envelope JSON que um
portão consegue ler.

\begin{center}
\begin{tabular}{ll}
\toprule
\textbf{Agente} & \textbf{A pergunta que ele responde} \\
\midrule
05 Gate estrutural & A missão tem estrutura construível? \\
08 Avaliação pedagógica & Cumpre currículo e protege o estudante? \\
09 Experiência do estudante & Alguém dessa idade consegue viver isso? \\
29 QA de entregáveis & A construção é fiel ao que foi especificado? \\
\bottomrule
\end{tabular}
\end{center}

\subsection{Redundância não é cobertura}

Três critérios do agente 09 eram testes idênticos a três critérios do agente 05 --- mesma
pergunta, mesmo modo de falha. Foram removidos. Knight e Leveson~\cite{knight1986} mostraram,
em 27 implementações independentes submetidas a um milhão de testes, que falhas correlacionadas
anulam o ganho da redundância. Quando os dois ``avaliadores independentes'' são prompts
diferentes do \emph{mesmo modelo}, sobre o \emph{mesmo texto}, a correlação é ainda maior.

\subsection{O contraponto honesto}

A literatura recente é cética quanto a sistemas multiagente: sob orçamento fixo de
\emph{tokens}, um agente único frequentemente iguala ou supera arquiteturas
multiagente~\cite{tran2026}. A condição de contorno declarada pelos próprios autores é que o
multiagente vence \textit{``quando a utilização efetiva de contexto de um agente único está
degradada''} --- que é exatamente o regime desta esteira. E Cemri et al.~\cite{cemri2025},
sobre mais de 1.600 execuções anotadas, mostram que decompor cria uma classe nova de falha:
contratos de interface. Por isso cada fronteira aqui tem envelope JSON verificado por teste.

\section{As etapas}

\begin{center}
\begin{tikzpicture}[node distance=0.42cm and 0.5cm]
\node[etapa, fill=fundoverde, draw=verde!60] (p1) {01--04\\ núcleo e one-pager};
\node[etapa, fill=fundoverde, draw=verde!60, right=of p1] (p2) {05--07\\ gate estrutural};
\node[etapa, fill=fundoverde, draw=verde!60, right=of p2] (p3) {08--11\\ avaliações e gate final};
\node[etapa, fill=fundoverde, draw=verde!60, right=of p3] (p4) {12--15\\ spec + currículo};
\node[etapa, right=of p4] (d1) {16--20\\ lacunas, ambiente, briefing};

\node[etapa, below=1.15cm of p1] (d2) {21--24\\ jornada, fluxo, copy, CAA};
\node[etapa, right=of d2] (d3) {25--28\\ telas e protótipo};
\node[etapa, right=of d3] (d4) {29--31\\ QA de design};
\node[etapa, fill=fundoambar, draw=ambar!70, right=of d4] (v1) {32--37\\ manifesto e handoff};
\node[etapa, fill=fundoambar, draw=ambar!70, right=of v1] (v2) {38--45\\ banca de 5 dimensões};

\node[etapa, fill=fundoambar, draw=ambar!70, below=1.15cm of d2] (v3) {46--48\\ PRD e one-pager};
\node[portao, right=of v3] (t1) {49 diagnóstico\\ (determinístico)};
\node[portao, right=of t1] (t2) {50 aprendizado};
\node[etapa, right=of t2, fill=white, draw=cinza] (fim) {PRD + one-pager\\ + diagnóstico};

\draw[seta] (p1) -- (p2); \draw[seta] (p2) -- (p3); \draw[seta] (p3) -- (p4); \draw[seta] (p4) -- (d1);
\draw[seta] (d1.south) to[out=-90,in=90] ($(d2.north)+(0,0.01)$);
\draw[seta] (d2) -- (d3); \draw[seta] (d3) -- (d4); \draw[seta] (d4) -- (v1); \draw[seta] (v1) -- (v2);
\draw[seta] (v2.south) to[out=-90,in=90] ($(v3.north)+(0,0.01)$);
\draw[seta] (v3) -- (t1); \draw[seta] (t1) -- (t2); \draw[seta] (t2) -- (fim);

\node[font=\scriptsize\bfseries, text=verde, above=0.05cm of p1, xshift=-0.1cm] {PRODUTO};
\node[font=\scriptsize\bfseries, text=azul, above=0.05cm of d1, xshift=0.1cm] {DESIGN};
\node[font=\scriptsize\bfseries, text=ambar, above=0.05cm of v1] {DEV};
\end{tikzpicture}
\captionof{figure}{Os 50 agentes, numerados na ordem de execução}
\end{center}

\subsection{Produto (01--15): a missão pode ser construída?}

Decide \emph{o que} a missão é. Começa pelo núcleo narrativo, passa por um micro-portão que
caça \emph{spoiler} e ausência de escolha significativa, completa o one-pager de 18 seções,
submete tudo aos 23 critérios do gate estrutural, avalia currículo e faixa etária em paralelo,
e fecha com a \textbf{spec} --- o contrato de 12 blocos contra o qual Design e Dev constroem.

\subsection{Design (16--31): como isso se constrói}

Fecha as lacunas da spec, decide o ambiente simulado, escreve o briefing de ofício, e então
constrói: jornada e fluxo em \textbf{diagramas Mermaid}, copy em duas passadas, camada de
acessibilidade, especificação de telas e um \textbf{protótipo HTML navegável} que abre com
duplo clique. Fecha com o QA de entregáveis, que audita fidelidade à spec.

\begin{quote}\small
\textbf{Nota sobre ferramentas.} A esteira não usa mais o Figma --- não há cliente oficial que o
sirva neste ambiente, e board em ferramenta que ninguém abre não é entregável. Jornada e fluxo
saem em Mermaid, as telas saem no Stitch e como especificação em Markdown, e o artefato visual
é o protótipo HTML.
\end{quote}

\subsection{Dev (32--48): isso é implementável?}

Confere a ficha do manifesto, produz o handoff e o brief de assets, valida o contrato de
telemetria, e submete tudo a uma banca de cinco dimensões em paralelo --- produto, design,
pedagogia, dados/LGPD e arquitetura. Consolida, escreve o PRD e, por último, o one-pager
pedagógico.

\subsection{Transversal (49--50): o que aconteceu}

O diagnóstico determinístico monta o dossiê a partir dos envelopes que os portões gravaram; o
agente de aprendizado escreve a leitura em prosa e registra o que a execução ensinou.

\section{Os seis ciclos de correção}

Quando um portão reprova, quem decide repetir é o \textbf{orquestrador}, em código --- não uma
instrução em linguagem natural dentro do prompt, que o modelo pode ou não seguir.

\begin{center}
\begin{tikzpicture}[node distance=0.55cm and 1.15cm]
\node[etapa] (a) {autor\\ produz};
\node[portao, right=of a] (g) {portão\\ julga};
\node[etapa, right=of g] (r) {corretor\\ conserta};
\node[etapa, right=1.9cm of r, fill=fundoverde, draw=verde!60] (s) {aprovado:\\ segue};
\draw[seta] (a) -- (g);
\draw[seta] (g) -- node[font=\scriptsize, above, text=vermelho] {reprova} (r);
\draw[seta] (g.north) to[out=45,in=135] node[font=\scriptsize, above] {aprova} (s.north);
\draw[setavolta] (r.south) to[out=-90,in=-90] node[font=\scriptsize, below, text=ambar] {reavalia (teto de rodadas)} (g.south);
\end{tikzpicture}
\captionof{figure}{O padrão comum aos seis ciclos: quem julga não conserta, quem conserta não julga}
\end{center}

\begin{center}
\begin{longtable}{p{2.5cm}p{4.3cm}p{7.6cm}}
\toprule
\textbf{Ciclo} & \textbf{Portão $\rightarrow$ corretor} & \textbf{O que resolve} \\
\midrule
\endfirsthead
Núcleo & micro-gate $\rightarrow$ 01 & \emph{Spoiler}, falta de segredo, escolha sem consequência. \\
\addlinespace
Gate estrutural & 05 $\rightarrow$ 07 revisor & Os 23 critérios de estrutura. O gate escreve o texto de correção pronto para colar. \\
\addlinespace
Spec & currículo + manifesto $\rightarrow$ 12 & Código da BNCC inexistente ou de outro ano; campo do manifesto fora do vocabulário fechado. Conserto mecânico. \\
\addlinespace
QA de design & 29 $\rightarrow$ dono do ticket & Cada achado tem dono declarado; só os endereçados rodam de novo, e o protótipo é refeito depois. \\
\addlinespace
Handoff & telemetria $\rightarrow$ 33 & Contrato de eventos, papéis, canal de sala. O mais barato: a correção vem escrita na mensagem do portão. \\
\addlinespace
Banca & banca $\rightarrow$ dono do achado & Inclui donos da \emph{etapa anterior}: um erro de jornada detectado em Dev volta para o agente de jornada. \\
\bottomrule
\end{longtable}
\end{center}

\noindent
\textbf{Dois freios.} Cada ciclo tem teto próprio, e existe um orçamento global de
reexecuções para a execução inteira (\cod{ORCAMENTO\_REEXECUCOES}, padrão 8). Estourado
qualquer um dos dois, a esteira para de tentar, marca pendência humana e segue até o fim.

\noindent
\textbf{O que não volta.} Achado cujo dono está na etapa de \emph{Produto} não reabre a spec:
reabri-la invalidaria tudo que o Design construiu sobre ela. Vira pendência nomeada no
diagnóstico, para decisão humana.

\section{Os portões, e o que fazer quando reprovam}

Há dois tipos. \textbf{Determinísticos} são funções Python: não gastam \emph{token}, não erram
e não mudam de opinião. \textbf{Avaliadores} são agentes que aplicam rubrica --- e aqui vale
saber que avaliadores baseados em modelos de linguagem têm vieses documentados de posição,
verbosidade e leniência~\cite{zheng2023}, razão pela qual toda regra que \emph{pode} ser
mecânica foi tirada deles.

\begin{center}
\begin{longtable}{p{3.5cm}p{1.6cm}p{9.3cm}}
\toprule
\textbf{Portão} & \textbf{Tipo} & \textbf{Reprova quando} \\
\midrule
\endfirsthead
Micro-gate (02) & rubrica & O briefing ou o núcleo entrega a virada; a escolha não muda nada. \\
Gate estrutural (05) & rubrica & Falta mecanismo, caminho sem tratamento, ambiente fora da lista. \\
Gate final (10) & rubrica & Os três relatórios avaliaram versões diferentes; correção ficou só no relatório. \\
Currículo BNCC (13) & \textbf{código} & Código inexistente, de outro ano, ou sem letramento digital. \\
Manifesto (14 e 32) & \textbf{código} & Campo do bloco 0 ausente ou fora do vocabulário fechado. \\
QA de design (29) & rubrica & Construção infiel à spec; diagrama ausente ou inválido. \\
Telemetria (35) & \textbf{código} & Evento fora do contrato real da API. \\
Rastro de lacunas (37) & \textbf{código} & Marcadores de lacuna abertos nos documentos. \\
Banca (38--44) & rubrica & Qualquer um dos critérios bloqueantes das cinco dimensões. \\
Pendências do PRD (47) & \textbf{código} & Marcadores que o redator deixou por falta de informação. \\
\bottomrule
\end{longtable}
\end{center}

\section{Como ler o diagnóstico}

O \cod{diagnostico-execucao.md} é o primeiro arquivo a abrir. Ele tem sempre a mesma ordem:

\begin{enumerate}[leftmargin=*,itemsep=1pt]
  \item \textbf{Veredito e contagem} --- impedimentos, ressalvas, decisões humanas.
  \item \textbf{Como cada portão se saiu} --- tabela com resultado e rodadas gastas.
  \item \textbf{O que reprovou, na ordem} --- a sequência dos fatos.
  \item \textbf{O que precisa ser consertado, por dono} --- agrupado por responsável.
  \item \textbf{O que só uma pessoa decide} --- o que nenhuma rodada resolve.
  \item \textbf{O que saiu provisório} --- os documentos que dependem de algo reprovado.
  \item \textbf{Documentos ausentes} --- sintoma de agente que não produziu nada.
  \item \textbf{Por onde começar} --- a ação que destrava mais coisa de uma vez.
\end{enumerate}

\noindent
Regra prática: \Vermelho{impedimento} é o que impede construir; \Ambar{ressalva} é risco
aceito e registrado; \textbf{decisão humana} é o que a esteira não tem autoridade para
resolver --- escolha de produto, licença de imagem, validação pedagógica.

\section{O identificador da missão}
\label{sec:id}

O código da experiência carrega ano e bimestre:

\begin{center}
\begin{tabular}{lll}
\toprule
\textbf{Código} & \textbf{Significa} & \textbf{Campos que têm de concordar} \\
\midrule
\cod{exp01ef1b} & 1.\textsuperscript{o} ano EF, 1.\textsuperscript{o} bimestre & \cod{series: ["1EF"]} · \cod{unidade: 1} \\
\cod{exp05ef3b} & 5.\textsuperscript{o} ano EF, 3.\textsuperscript{o} bimestre & \cod{series: ["5EF"]} · \cod{unidade: 3} \\
\cod{exp02em2b} & 2.\textsuperscript{o} ano EM, 2.\textsuperscript{o} bimestre & \cod{series: ["2EM"]} · \cod{unidade: 2} \\
\bottomrule
\end{tabular}
\end{center}

\noindent
Um portão determinístico confere que \cod{series} e \cod{unidade} concordam com o
identificador --- divergência trava, porque é ele que nomeia tabela, pasta e \emph{branch}.
Missão transversal (sem série) e ferramenta de professor não têm ano nem bimestre: seguem no
formato antigo \cod{exp-NN-<slug>}.

\section{Configuração}

\begin{center}
\begin{tabular}{p{5.1cm}p{1.5cm}p{8.1cm}}
\toprule
\textbf{Variável} & \textbf{Padrão} & \textbf{O que faz} \\
\midrule
\cod{GATE\_ESTRUTURAL\_MAX\_RODADAS} & 3 & Rodadas do ciclo gate $\leftrightarrow$ revisor. \\
\cod{SPEC\_MAX\_RODADAS} & 2 & Ciclo spec $\leftrightarrow$ currículo e manifesto. \\
\cod{QA\_DESIGN\_MAX\_RODADAS} & 2 & Ciclo QA $\leftrightarrow$ retrabalho de design. \\
\cod{HANDOFF\_MAX\_RODADAS} & 2 & Ciclo handoff $\leftrightarrow$ telemetria. \\
\cod{BANCA\_MAX\_RODADAS} & 2 & Ciclo banca $\leftrightarrow$ dono do achado. \\
\cod{ORCAMENTO\_REEXECUCOES} & 8 & Teto global de agentes reexecutados. \\
\cod{ESTEIRA\_ESTRITA} & true & Reprovação marca herança degradada e o PRD sai 🔴; desligado, vira ressalva. \\
\cod{PIPELINE\_STAGES} & todas & Permite rodar só uma etapa, sobre um estado salvo. \\
\bottomrule
\end{tabular}
\end{center}

\noindent
Qualquer teto em 1 desliga o ciclo correspondente: a esteira avalia e reporta, sem corrigir.

\section{Três recomendações de uso}

\begin{description}[leftmargin=1.1cm,style=nextline,itemsep=3pt]
  \item[Comece pelo diagnóstico, não pelo PRD.] Ele diz em trinta segundos se vale ler o resto.
  \item[Trate \cod{[PROVISÓRIO]} como rascunho, não como entrega.] O marcador existe porque o
  documento foi construído sobre algo que um portão reprovou.
  \item[Rode uma etapa isolada quando estiver iterando.] \cod{PIPELINE\_STAGES} sobre um estado
  salvo custa minutos em vez de uma execução inteira, e a entrada fica fixa entre tentativas.
\end{description}

\begin{thebibliography}{99}

\bibitem{cemri2025}
CEMRI, M.; PAN, M. Z.; YANG, S. et al. \textit{Why Do Multi-Agent LLM Systems Fail?}
NeurIPS 2025, Datasets and Benchmarks Track. arXiv:2503.13657.

\bibitem{cne2022}
BRASIL. Conselho Nacional de Educação. \textit{Resolução CNE/CEB n.\textsuperscript{o} 1, de 4
de outubro de 2022}: normas sobre Computação na Educação Básica --- Complemento à BNCC.
Brasília: MEC, 2022.

\bibitem{cowan2001}
COWAN, N. The magical number 4 in short-term memory. \textit{Behavioral and Brain Sciences},
v. 24, n. 1, p. 87--114, 2001.

\bibitem{ifscale2025}
JAROSLAWICZ, D.; WHITING, B.; SHAH, P.; MAAMARI, K. \textit{How Many Instructions Can LLMs
Follow at Once?} arXiv:2507.11538, 2025.

\bibitem{jonsson2007}
JONSSON, A.; SVINGBY, G. The use of scoring rubrics: reliability, validity and educational
consequences. \textit{Educational Research Review}, v. 2, n. 2, p. 130--144, 2007.

\bibitem{knight1986}
KNIGHT, J. C.; LEVESON, N. G. An experimental evaluation of the assumption of independence in
multiversion programming. \textit{IEEE TSE}, v. SE-12, n. 1, p. 96--109, 1986.

\bibitem{liu2024}
LIU, N. F.; LIN, K.; HEWITT, J. et al. Lost in the middle: how language models use long
contexts. \textit{TACL}, v. 12, p. 157--173, 2024.

\bibitem{parnas1972}
PARNAS, D. L. On the criteria to be used in decomposing systems into modules.
\textit{Communications of the ACM}, v. 15, n. 12, p. 1053--1058, 1972.

\bibitem{sweller1988}
SWELLER, J. Cognitive load during problem solving: effects on learning.
\textit{Cognitive Science}, v. 12, n. 2, p. 257--285, 1988.

\bibitem{tran2026}
TRAN, D.; KIELA, D. \textit{Single-Agent LLMs Outperform Multi-Agent Systems on Multi-Hop
Reasoning Under Equal Thinking Token Budgets}. arXiv:2604.02460, 2026.

\bibitem{zheng2023}
ZHENG, L.; CHIANG, W.-L.; SHENG, Y. et al. Judging LLM-as-a-judge with MT-Bench and Chatbot
Arena. NeurIPS 2023. arXiv:2306.05685.

\end{thebibliography}

\end{document}
